% RESEARCH_DRAFT_SHA256: e24ef756503818a78300cc1176469dd2c1f060d9c6d9c7daf8593ce9ede0f12d
% Layout and source-reference wrappers expanded for the web renderer.
\documentclass[11pt]{article}
\usepackage[margin=25mm]{geometry}
\usepackage{fontspec,unicode-math}
\setmainfont{Latin Modern Roman}
\setmathfont{Latin Modern Math}
\usepackage{amsmath,amssymb,mathtools,xurl,hyperref}
\hypersetup{colorlinks=true,urlcolor=blue,linkcolor=blue}
\setcounter{secnumdepth}{0}
\setlength{\parindent}{0pt}
\setlength{\parskip}{5pt}
\setlength{\emergencystretch}{3em}
\newcommand{\sref}[2]{\hyperlink{source:#1:#2}{[S#2]}}
\newcommand{\cataloguescope}[1]{\paragraph{Source record.}#1}
\newcommand{\floor}[1]{\left\lfloor#1\right\rfloor}
\begin{document}
\section*{Research-attempt provenance}
This standalone research attempt retains the complete original definition body below,
including its historical source and preparation notes. Those notes describe the
earlier catalogue preparation; the new attempt and its verification are identified
separately after that body. No original ID, catalogue status, or source statement is changed.

% UnsolvedMath ID 20001227; source catalogue code AIM-COMPUTATION-0065
\section{Gaussian conditional independence as an algebraic minor}\label{20001227}
{\small UnsolvedMath ID 20001227 · Computer Science}
\subsection{Definitions and mathematical statement}

\paragraph{Shared notation.}Standard number systems are $\mathbb N,\mathbb Z,\mathbb Q,\mathbb R,\mathbb C$. Specialized notation and index ranges follow the cited source and the definitions below.


Continuous random variables take values in spaces such as $\mathbb R^d$. A multivariate Gaussian model is parameterized by a mean and covariance matrix. In Gaussian models, conditional independence can be expressed by vanishing covariance or concentration minors under the appropriate nonsingularity assumptions. Microarray measurements are an application requiring a noise model and inferential target.
\paragraph{Record scope and qualifications.}
The catalogue title is retained as metadata. Named certificates, contracts, bounds or frameworks in that title are not asserted as results of the original source extract.
The source mentions Gaussian assumptions as a possible direction, not a restriction imposed on every continuous model. No specific microarray analysis or conditional-independence test is prescribed. The supplied catalogue marks this record partially solved; its present research status was not independently checked.
The source text below is rendered literally where necessary; damaged formulas and composite extracts are not silently repaired into a new mathematical claim.
\paragraph{Statement recorded in the supplied source.}
Question 50 (Kuhnt). What about adding continuous random variables into these prob-lems? Might an algebraic framework be useful? Dinwoodie: Can we try to apply algebraic techniques to microarray data? Meek: Gaussian (or similar) assumptions on the random variables might be possible to handle with algebraic techniques. \href{https://aimath.org/WWN/compalgstat/compalgstat.pdf}{[S2]}
\paragraph{Source status.}
The supplied catalogue labels this source record partially solved; that metadata was checked against the supplied snapshot on 2026-10-11. The label is preserved for provenance and is not an independent verification of present open status or mathematical correctness.
\subsection{Short English statement}
The source asks whether algebraic methods can handle continuous variables and gene-expression data.
\subsection{Sources}
\begin{itemize}
\item[S1] Ulam AI and the UnsolvedMath Contributors, UnsolvedMath: A Curated Collection of Open Mathematics Problems, record 20001227 (AIM-COMPUTATION-0065). Catalogue curation and metadata: CC BY 4.0. The source extract is retained; definitions, English summaries and qualifications are editorial additions.
\url{https://huggingface.co/datasets/ulamai/UnsolvedMath}.
\item[S2] AIM Workshop Problem Lists, source linked by catalogue record AIM-COMPUTATION-0065.
\url{https://aimath.org/WWN/compalgstat/compalgstat.pdf}.
\end{itemize}
\paragraph{Preparation and review.}
The original source and UnsolvedMath contributors retain attribution for the source extract. Codex (GPT-6; exact serving revision not exposed) assisted with definitions, English summaries, scope notes and formatting. No new human mathematical authorship or independently verified solution is asserted. Permanent source ID 20001227 and code AIM-COMPUTATION-0065 checked against the complete supplied registry and current repository. Distinct upstream records, including repeated questions and broad or damaged extracts, are retained for complete coverage of the source CS category. A repeated record is not a new independent conjecture. Complete mathematical-equivalence and current-literature review has not been performed. Preparation and final review follow contribution prompts C and D, with source gaps disclosed. Structural and rendering checks do not establish a new theorem or open status.
\label{end:20001227}

\section{Literature and context check}
AIM Question 50 mentions Gaussian assumptions as a possibility, without restricting all continuous models to that class. The calculation below reproves a Gaussian criterion and tests what fails without Gaussianity or nonsingularity. It contains no microarray-data analysis or novelty claim.
\section{Research attempt: A Gaussian algebraic criterion and its two boundary failures}
Let $(X_i,X_j,Z)$ be a centered jointly Gaussian vector with positive definite covariance, and write $W=(X_i,X_j)^T$, $C=\operatorname{Cov}(Z)$, $D=\operatorname{Cov}(W,Z)$ and $B=\operatorname{Cov}(W)$. Put $R=W-DC^{-1}Z$. Its covariance with $Z$ is zero and its covariance is $S=B-DC^{-1}D^T$. This matrix is positive definite: for nonzero $u$, the variance of $u^TW-u^TDC^{-1}Z$ is positive by positive definiteness of the original covariance.

The Gaussian vector $(R,Z)$ has block diagonal covariance. Its density factors because the inverse covariance is block diagonal and its determinant is $\det S\det C$; the usual Gaussian density formula therefore gives a product of its two marginal densities. Thus $R$ and $Z$ are independent. Conditional on $Z=z$, $W$ is a translate of the Gaussian $R$, so conditional independence of its two scalar coordinates is equivalent to independence of the coordinates of $R$. Those are independent when $S_{12}=0$, by the same density factorization; conversely independence makes their covariance zero. Consequently
\[
X_i\mathbin{\perp\!\!\!\perp}X_j\mid Z
\quad\Longleftrightarrow\quad
\Sigma_{ij}\det C-\Sigma_{iZ}\operatorname{adj}(C)\Sigma_{Zj}=0.
\]
Here $C^{-1}=\operatorname{adj}(C)/\det C$ follows from the cofactor identity $C\operatorname{adj}(C)=(\det C)I$, and $\det C>0$. This proves the selected polynomial criterion, including its domain.

The first attempted extension is to general continuous variables. Take $X$ uniform on $[-1,1]$, $E$ an independent centered Gaussian with positive variance, and $Y=X^2+E$. Their joint distribution has a density, and
\[
\operatorname{Cov}(X,Y)=E X^3-E X\,E X^2=0.
\]
They are dependent: $E[Y\mid X]=X^2$ is nonconstant, whereas independence and integrability would force this conditional expectation to equal the constant $E Y$. Thus a vanishing covariance polynomial does not identify independence for arbitrary continuous distributions.

The second attempted extension clears denominators and includes singular boundary covariances without reanalysis. Take $Z=0$ deterministically and $X_i=X_j=G$ for a nondegenerate centered Gaussian $G$ of variance one. The displayed polynomial is zero: $\det C=0$ and both cross-covariances with $Z$ vanish. But $X_i$ and $X_j$ are not conditionally independent given the constant $Z$; their covariance is one. The polynomial equation alone has acquired spurious boundary cases.
\section{Verification and adversarial audit}
Gaussian sufficiency uses an explicit density factorization, and necessity uses covariance of independent integrable variables. The non-Gaussian example has a joint density despite its nonlinear conditional mean. The singular example is deliberately outside the positive-definite hypothesis; it refutes only the proposed hypothesis-free extension.
\paragraph{Audit corrections.} Positive definiteness is retained through denominator clearing, and the Gaussian special case is not silently substituted for every continuous distribution in the source.
\paragraph{External computation.} None; symbolic arguments and boundary checks only.
\section{Scope and remaining gaps}
This is a standard Gaussian special case with proved obstructions to two extensions. The broad continuous-model and microarray agenda remains unspecified and is not solved; no general conditional-independence test or statistical calibration is supplied.
\paragraph{Integration review.} Codex (GPT-6; exact serving revision not exposed) reviewed the written arguments and expanded layout and source-reference wrappers. Repository metadata and the production web renderer were checked using Python 3.13 and JavaScript with MathJax 3.2.2. This model-assisted review is not independent human proof verification; no PDF compilation or formal proof checking is asserted.

\section{Completion Estimate}
\noindent\textbf{COMPLETION: 0\%}
The calculations below the preserved source statement investigate a concrete route,
but establish no substantiated new progress toward the entire original target.
Known special cases and failed sufficient conditions are not credited as a solution.
\section{References}
The original source references and catalogue attribution are preserved in the
Sources section above. The following URLs were actually consulted for this attempt;
retention of an additional original reference does not claim fresh consultation of it.
\begin{itemize}
\item \url{https://aimath.org/WWN/compalgstat/compalgstat.pdf} (consulted 2026-10-11). 
\end{itemize}
\section{Model and contributor disclosure}
OpenAI gpt-6-astra, requested reasoning configuration ultra, generated this research
attempt on 2026-10-11. The exact serving revision is not exposed. Codex GPT-6
performed orchestration, repository-format review, and a model-assisted mathematical
audit of the written arguments; its exact serving revision
is also not exposed. The mathematical route, boundary checks and final adversarial
self-audit followed contribution prompts A then B in
\texttt{docs/contribution\_prompts.md}, after reading \texttt{CONTRIBUTING.md}.
Python 3.14.2 assembled the standalone documents; any mathematical computation is
specified in Verification. No theorem prover was used. No human mathematical
authorship, substantial human proof correction, or independent human proof checking
is asserted. Model-assisted review, source retrieval, and structural checks do not
establish correctness or novelty. Full private reasoning transcripts are unavailable
as reproducibility artifacts; the complete arguments, inputs and outcomes actually
used are stated here or in the companion batch audit.
\section{Final status}
Unresolved as a new research contribution. This is a concrete route investigation,
not a claim of new complete resolution or a current-status certification. A known
answer reproduced for context does not become a new contribution. The catalogue's
source-status label is preserved; this attempt status does not assert that a
source-labelled solved question is open.
\end{document}
